\section{Описание предполагаемого метода}

\subsection{Подготовка наблюдений}

Сначала формируется метрическая модель объекта: координаты опорных точек, принадлежность граней и, при наличии маркеров, соответствие их идентификаторов геометрии. Затем на изображениях обнаруживаются видимые элементы, уточняются координаты углов и отбрасываются неоднозначные соответствия. Для каждого набора сохраняются идентификаторы камеры, положения объекта и точки. Пиксельные координаты и модель проекции должны соответствовать одному представлению изображения: нельзя одновременно использовать исправленные изображения и повторно применять к ним исходную дисторсию.

Наблюдения желательно получать при нескольких положениях и ориентациях мишени. Повторение практически одинакового ракурса увеличивает число отсчётов, но не заменяет геометрического разнообразия. Отдельно проверяются достаточность числа точек, их распределение по изображению и видимость объекта перед камерой.

\subsection{Начальная оценка внешних параметров}

Для каждой допустимой пары «камера --- положение объекта» решается задача PnP. Полученная матрица $A_{is}$ переводит координаты объекта в систему камеры, то есть в идеальной модели
\begin{equation}
  A_{is}=T_iP_s.
\end{equation}
EPnP~\cite{lepetit2009} может использоваться для начального решения и поддерживает плоские и пространственные конфигурации точек. При наблюдении только одной грани необходимо отдельно контролировать неоднозначность и устойчивость оценки, проверяя положительность глубин и согласованность с другими наблюдениями. Оценка дополняется устойчивой отбраковкой выбросов и локальным уточнением ошибки проекции.

Если камеры $i$ и $j$ наблюдают одну позу объекта, начальное относительное преобразование вычисляется как
\begin{equation}
  T_{j\leftarrow i}^{(s)}=A_{js}A_{is}^{-1}.
  \label{eq:initial}
\end{equation}
При отсутствии общих наблюдений у пары камер связь может быть восстановлена по пути в связном графе. Полученные оценки служат инициализацией; итоговое решение не должно зависеть только от одного такого пути. Для проверки соглашений о координатах используется обратное проецирование известных точек по формуле~\eqref{eq:projection}.

\subsection{Совместное уточнение по изображениям}

После инициализации внешние параметры и позы мишени уточняются одновременно. Для множества корректных наблюдений $\Omega_I$ определяется функция
\begin{equation}
  E_I=\sum_{(i,j,s)\in\Omega_I}
  \rho_I\!\left(
    \frac{\|\mathbf u_{ijs}-\widehat{\mathbf u}_{ijs}\|^2}
         {\sigma_{u,ijs}^2}
  \right),
  \label{eq:imagecost}
\end{equation}
где $\sigma_{u,ijs}$ характеризует масштаб ошибки локализации, а $\rho_I$ --- робастная функция потерь. При отсутствии индивидуальных оценок точности допускается общий масштаб, который должен быть явно указан. Оптимизация проводится по $T_i$ и $P_s$ с фиксированным $T_0=I$; вращения параметризуются с сохранением принадлежности $SO(3)$. Внутренние параметры и геометрия объекта в базовой постановке остаются фиксированными.

Такая задача является специализированным вариантом совместной оптимизации~\cite{triggs2000}. Робастная функция снижает влияние отдельных неверных наблюдений, но не устраняет систематическую ошибку геометрии или внутренней калибровки. Поэтому после оптимизации требуется анализ остатков по камерам, положениям и участкам изображения.

\subsection{Дополнение измерениями глубины}

Если доступны надёжные карты глубины, можно использовать точки внутренних областей граней. Для каждого пикселя сначала восстанавливаются координаты в системе датчика глубины. Если значение $z$ означает осевую координату, то для исправленного пикселя $\mathbf v=(u,v,1)^\mathsf T$
\begin{equation}
  \mathbf q=zK_D^{-1}\mathbf v.
\end{equation}
Если же значение $r$ означает расстояние вдоль луча, используется другая формула:
\begin{equation}
  \mathbf q=r\frac{K_D^{-1}\mathbf v}
                       {\|K_D^{-1}\mathbf v\|}.
\end{equation}
Здесь $K_D$ относится к каналу глубины. Перед использованием точку необходимо перевести в выбранную систему $C_i$ по известному преобразованию между каналами. Простое совпадение размеров изображений не гарантирует их геометрического совмещения.

Пусть грань $f$ в системе объекта описывается единичной нормалью $\mathbf n_f$ и уравнением $\mathbf n_f^\mathsf T\mathbf x+b_f=0$. Для измеренной точки $\mathbf q_{iks}$, уже выраженной в $C_i$, остаток равен
\begin{equation}
  r^D_{iks}=\mathbf n_f^\mathsf T
    \operatorname{xyz}\!\left((T_iP_s)^{-1}
      \widetilde{\mathbf q}_{iks}\right)+b_f,
\end{equation}
где $f$ --- установленная принадлежность точки грани, а $\operatorname{xyz}$ извлекает три пространственные координаты. Наборы точек ограничиваются внутренними областями граней, чтобы уменьшить влияние смешанных измерений на границах.

Совместный критерий записывается как
\begin{align}
  E &= E_I+\lambda E_D,\\
  E_D &=\sum_{(i,k,s)\in\Omega_D}
    \rho_D\!\left(\frac{(r^D_{iks})^2}{\sigma_{d,iks}^2}\right),
\end{align}
где $\sigma_{d,iks}$ --- масштаб погрешности расстояния до плоскости, а $\lambda$ управляет вкладом глубины. Этот масштаб должен учитывать распространение погрешности измерения вдоль луча на нормаль грани. Нормирование делает остатки безразмерными и позволяет согласовать их вклад. Влияние числа и пространственной корреляции точек также следует учитывать при выборе выборки и веса. Значение $\lambda=0$ соответствует базовой калибровке по изображениям.

Измерения одной плоскости не ограничивают одинаково все компоненты позы. Поэтому глубина дополняет визуальные соответствия, а достаточность геометрии оценивается для набора наблюдений в целом. Конкретная модель достоверности и настройки робастной функции должны выбираться после изучения характеристик камеры, а не назначаться как универсальные константы.

\subsection{Критерии последующей проверки}

Сопоставление базового и расширенного вариантов предполагается проводить на одинаковых наборах наблюдений. Следует анализировать ошибку репроекции на отложенных положениях мишени, устойчивость оценок при изменении состава кадров и согласованность независимо оценённых парных преобразований по циклам графа до совместной оптимизации. Последняя является внутренней проверкой и сама по себе не доказывает абсолютную точность. Циклы, составленные только из итоговых матриц $T_jT_i^{-1}$, замыкаются тождественно и не являются критерием качества.

При проверке на отложенных положениях внешние параметры камер фиксируются. Для каждого положения оценивается одна общая поза объекта по части наблюдений, а ошибка вычисляется на оставшейся части. Независимая подгонка позы для каждой камеры не проверяла бы их взаимную калибровку.

При наличии независимого эталона можно измерять ошибки вращения и переноса. Без такого эталона допустимо сообщать ошибки проекции, геометрические остатки и повторяемость, но нельзя отождествлять их с истинной погрешностью внешних параметров. Проверка на глубине должна по возможности использовать данные, не участвовавшие в оптимизации.

Отдельный предполагаемый анализ --- внесение контролируемых возмущений в размеры мишени, углы между гранями и координаты меток. Он позволит оценить чувствительность решения к неточности изготовления и сопоставить её с влиянием ошибок изображения и глубины. Такой анализ пока является планом исследования; утверждений о достигнутой точности и преимуществе конкретной конфигурации в настоящей работе не делается.
